\section*{الفصل \ref{ch:b2:ellingham} --- مخططات إلينغهام والتعدين}

\begin{solution}{exo:b2:ellingham:1}
\ce{4Cu + O2 -> 2Cu2O}؛ \ce{4/3Al + O2 -> 2/3Al2O3}؛ \ce{2C + O2 -> 2CO}؛
\ce{C + O2 -> CO2}.
\end{solution}

\begin{solution}{exo:b2:ellingham:2}
$\Delta_r H^\circ = 2(-350.46) = \qty{-700.92}{kJ}$؛ $\Delta_r S^\circ = 2(43.65)
- 2(41.63) - 205.15 = \qty{-201.1}{J/K}$. ومنه $\Delta_r G^\circ = -700.9 +
0.2011\,T$ (\unit{kJ}، و$T$ بوحدة \unit{K})، وهي صالحة حتى نقطة انصهار الزنك.
\end{solution}

\begin{solution}{exo:b2:ellingham:3}
عند \qty{1500}{K} يقع خط \ce{Cr2O3} قرب \qty{-495}{kJ}. فالفلزات التي تقع تحته، ومن ثم
تستطيع اختزال أكسيد الكروم: السيليكون، والتيتانيوم، والألومنيوم، والمغنيسيوم،
والكالسيوم. والتي تقع فوقه، فلا تستطيع: النحاس، والحديد، والزنك. أما الكربون، عند \qty{-488}{kJ}، فيقع
فوقه بقليل: فلا يختزل أكسيد الكروم إلا عند درجة حرارة أعلى بقليل،
ويعطي عندئذ كرومًا محمَّلًا بالكربيد --- وهذا أحد أسباب صنع الكروم المعدّ
لسبائك خالية من الكربون \omterm{def:b2:ellingham:aluminothermy}{بالاختزال الألومنيوحراري}.
\end{solution}

\begin{solution}{exo:b2:ellingham:4}
الميل هو $-\Delta_r S^\circ$، ويحكمه تغير كمية الغاز.
\ce{C + O2 -> CO2}: مول واحد من الغاز يعطي مولًا واحدًا، $\Delta_r S^\circ \approx 0$،
خط أفقي. \ce{2C + O2 -> 2CO}: مول يعطي مولين، $\Delta_r S^\circ > 0$، فينحدر
الخط. \ce{2CO + O2 -> 2CO2}: ثلاثة تعطي اثنين، $\Delta_r S^\circ < 0$، فيصعد
الخط، مثل خطوط الفلزات.
\end{solution}

\begin{solution}{exo:b2:ellingham:5}
$\Delta_r G^\circ(\qty{600}{K}) = -181.6 + 600 \times 0.2165 = \qty{-51.7}{kJ}$،
$p_{\ce{O2}} = p^\circ\exp(-51\,710/(8.314 \times 600)) = \qty{3.2e-5}{bar}$. والهواء
فيه \qty{0.21}{bar}، أي أكثر بكثير: فيتأكسد الزئبق عند \qty{600}{K} (ببطء
--- وهكذا حُضّر الأكسيد أول مرة)، ولا يتفكك الأكسيد إلا
فوق نحو \qty{730}{K}.
\end{solution}

\begin{solution}{exo:b2:ellingham:6}
$\Delta_r G^\circ = -1582.3 - (-743.5) = \qty{-838.8}{kJ}$. فالكاشفان
جسمان صلبان لا يتماسّان إلا عند سطوح حبيباتهما، وللتفاعل
طاقة تنشيط عالية: فيجب إشعاله موضعيًا (بفتيل من المغنيسيوم)؛
ثم تُبقيه الحرارة التي يحررها مستمرًا.
\end{solution}

\begin{solution}{exo:b2:ellingham:7}
يقطع خط \ce{2C + O2 -> 2CO} خط \ce{2Fe + O2 -> 2FeO} قرب
\qty{1040}{K}؛ وفوق ذلك يكون لتفاعل \ce{FeO + C -> Fe + CO} \omterm{def:b2:reaction-free-energy:reaction-gibbs}{طاقة جيبس قياسية}
سالبة.
\end{solution}

\begin{solution}{exo:b2:ellingham:8}
$\Delta_r G^\circ = 2(-209.1) - (-396.0) = \qty{-22.2}{kJ}$، $K^\circ =
\exp(22\,200/(8.314 \times 1100)) = 11.3$. ومنه $x^2/(1 - x) = 11.3$، $x =
0.92$. وعند التبريد إلى \qty{700}{K} يعود التوازن نحو \ce{CO2}
($x = 0.016$): فيتفاعل أحادي أكسيد الكربون تفاعل عدم تناسب، \ce{2CO -> C + CO2}، ويترسب
السخام --- ببطء، ما لم يحفّزه سطح كالحديد.
\end{solution}

\begin{solution}{exo:b2:ellingham:9}
يصعد خط الماء (ثلاثة مولات من الغاز تعطي اثنين) أقل حدة من
خط $\ce{CO}/\ce{CO2}$، الذي هو أيضًا ثلاثة مقابل اثنين لكن بخسارة إنتروبيا أكبر. ويتقاطع
الخطان قرب \qty{1100}{K}: تحته يكون خط $\ce{CO}/\ce{CO2}$ أدنى ويكون
أحادي أكسيد الكربون المختزل الأقوى؛ وفوقه يكون خط الماء أدنى ويتفوق
الهيدروجين. وهو القول نفسه في انزياح توازن غاز
الماء \ce{CO + H2O <=> CO2 + H2} نحو اليسار عند درجة حرارة مرتفعة.
\end{solution}

\begin{solution}{exo:b2:ellingham:10}
لكل مول من \ce{O2}، أي من أجل \ce{2MgO + Si -> SiO2 + 2Mg(g)}: $\Delta_r
G^\circ = -643.7 - (-845.5) = \qty{+201.8}{kJ}$. ومع سيليكون وسيليكا
نشاطية كل منهما 1، $\Delta_r G = \Delta_r G^\circ + RT\ln(p_{\ce{Mg}}/p^\circ)^2$، وهي سالبة
حين $p_{\ce{Mg}} < p^\circ\exp(-201\,800/(2 \times 8.314 \times 1500)) =
\qty{3e-4}{bar}$. ففرن التفريغ يضخ بخار المغنيسيوم بسرعة تكوّنه نفسها
ويكثفه على سطح بارد، فلا يبلغ ضغطه أبدًا
تلك القيمة؛ والجير المضاف إلى الشحنة يربط السيليكا ويخفض نشاطيتها،
وهو ما يساعد أكثر.
\end{solution}

\begin{solution}{exo:b2:ellingham:11}
أربعة أنواع، وتفاعل واحد، وثلاثة أطوار (جسمان صلبان والغاز): $v = 4 - 1 +
2 - 3 = 2$. فعند درجة حرارة وضغط معلومين تتحدد النسبة $p_{\ce{CO}}/p_{\ce{CO2}}$:
فلا يستطيع الغاز أن يعطي أحادي أكسيد الكربون كله لأكسيد الحديد،
ويبقى غاز القمة محتويًا على كثير منه --- يُحرق لتسخين
الهواء مسبقًا.
\end{solution}

\begin{solution}{exo:b2:ellingham:12}
$\log K = 2(0.34 + 0.41)/0.059 = 25.4$، $K = \num{3e25}$. ومع
$[\ce{Fe^2+}] = \qty{1.0}{mol/L}$، يكون $[\ce{Cu^2+}] = 1/K \approx
\qty{4e-26}{mol/L}$: فيُزال النحاس تمامًا.
\end{solution}

\begin{solution}{pb:b2:ellingham:1}
\textbf{1.} $\Delta_r H^\circ = \qty{-700.92}{kJ}$، $\Delta_r S^\circ =
\qty{-201.1}{J/K}$.
\textbf{2.} $\Delta_r G^\circ = -700.9 + 0.2011\,T$ (\unit{kJ})، تحت
\qty{692.7}{K}.
\textbf{3.} $\Delta_{\text{انصهار}}S = 7320/692.7 = \qty{10.6}{J/(K.mol)}$؛
$\Delta_{\text{تبخر}}S = 115\,300/1180.2 = \qty{97.7}{J/(K.mol)}$.
\textbf{4.} الزنك السائل: $\Delta_r H^\circ = -700.92 - 2(7.32) = \qty{-715.6}{kJ}$،
$\Delta_r S^\circ = -201.1 - 2(10.6) = \qty{-222.2}{J/K}$: $\Delta_r G^\circ =
-715.6 + 0.2222\,T$.
\textbf{5.} الزنك الغازي: $\Delta_r H^\circ = -715.6 - 2(115.3) = \qty{-946.2}{kJ}$،
$\Delta_r S^\circ = -222.2 - 2(97.7) = \qty{-417.7}{J/K}$: $\Delta_r G^\circ =
-946.2 + 0.4177\,T$.
\textbf{6.} الميول $0.201$ و$0.222$ و$0.418$ \unit{kJ/K}. فوق نقطة
الغليان يستهلك التفاعل ثلاثة مولات من الغاز بدل مول واحد: فتتضاعف خسارة
الإنتروبيا، ومن ثم الميل.
\textbf{7.} عند \qty{692.7}{K}: $-700.9 + 139.3 = -715.6 + 153.9 = \qty{-561.6}{kJ}$؛
وعند \qty{1180.2}{K}: $-715.6 + 262.3 = -946.2 + 492.9 = \qty{-453.3}{kJ}$. فتلتقي
الخطوط، كما ينبغي.
\textbf{8.} الميل $(-453.0 + 418.2)/200 = \qty{-0.174}{kJ/K}$، $a = -418.2 +
0.174 \times 1100 = \qty{-226.8}{kJ}$: $\Delta_r G^\circ = -226.8 - 0.174\,T$.
\textbf{9.} $\Delta_r S^\circ = \qty{+174}{J/K}$: مول واحد من الغاز يعطي مولين.
\textbf{10.} \ce{2ZnO + 2C -> 2Zn(g) + 2CO}، $\Delta_r G^\circ = (-226.8 - 0.174\,T)
- (-946.2 + 0.4177\,T) = 719.4 - 0.592\,T$ (\unit{kJ}).
\textbf{11.} تنعدم عند $T = 719.4/0.592 = \qty{1215}{K}$.
\textbf{12.} فوق نقطة غليانه يتكوّن الزنك بخارًا: فالتفاعل يصنع
غازًا من أجسام صلبة، وإنتروبياه كبيرة وموجبة، ويغادر البخار
الشحنة، وهو ما يفصل الفلز عن الخام وفحم الكوك بالتقطير.
\textbf{13.} أربعة أنواع، وتفاعل واحد، وشرط خاص واحد ($p_{\ce{Zn}} =
p_{\ce{CO}}$)، وثلاثة أطوار: $v = 4 - 1 - 1 + 2 - 3 = 1$. فتثبيت الضغط
يحدد درجة حرارة التوازن.
\textbf{14.} مول من بخار الزنك لكل مول من أحادي أكسيد الكربون:
$p_{\ce{Zn}} = p_{\ce{CO}} = \qty{0.5}{bar}$.
\textbf{15.} لكل مول من \ce{ZnO}، $\Delta_r G^\circ = \frac12(719.4 - 0.592 \times
1300) = \qty{-25.1}{kJ/mol}$، $K^\circ = \exp(25\,100/(8.314 \times 1300)) = 10$؛
$Q = 0.25 < K^\circ$: فيجري الاختزال.
\textbf{16.} فوق التقاطع بقليل يكون الاختزال مفضَّلًا منذ الآن (مع
\qty{0.5}{bar} من كل غاز، ابتداءً من نحو \qty{1170}{K})؛ والتسخين الأشد يكلّف وقودًا،
ويُبلي المعوجات الطينية، ولا يصنع زنكًا أكثر، لأن التفاعل يبلغ
نهايته على أي حال.
\textbf{17.} \ce{Zn(g) + CO2 -> ZnO + CO}.
\textbf{18.} يقع خط أكسيد الزنك تحت خط $\ce{CO}/\ce{CO2}$
($-445$ مقابل \qty{-357}{kJ} عند \qty{1200}{K}): فيختزل بخار الزنك ثنائي أكسيد
الكربون ويعود أكسيدًا. وفي المعوجة يُفني الكربون الساخن كل
\ce{CO2} (بودوار)؛ أما في المكثِّف الأبرد فلا شيء يفعل ذلك، و\ce{CO2} --- الآتي من
الهواء المتسرب، أو من \ce{2CO -> C + CO2} عند التبريد --- يُفسد الزنك فيحوّله إلى
مسحوق رمادي متأكسد.
\textbf{19.} التبريد البطيء يُبقي البخار مدة طويلة في المجال الذي
يعود فيه إلى التأكسد، ويعطي غبارًا دقيقًا بدل فلز سائل؛ والتكثيف
السريع إلى السائل يحدّ من الأمرين.
\textbf{20.} $10^6/65.38 = \qty{1.53e4}{mol}$ من الزنك، ومثلها من الكربون:
$1.53 \times 10^4 \times 12.011 = \qty{184}{kg}$ (وأكثر بكثير عمليًا، لتسخين
المعوجات).
\textbf{21.} $\Delta_r H^\circ = \frac12(-226.8 + 946.2) = \qty{+360}{kJ}$ لكل
مول من الزنك؛ ولكل طن، $1.53 \times 10^4 \times 360 = \qty{5.5e6}{kJ}$، أي
\qty{5.5}{GJ} --- يوفرها حرق مزيد من الفحم حول المعوجات.
\textbf{22.} يختزل الكربون أكسيد الزنك في الشروط القياسية فوق
$\boldsymbol{T \approx \qty{1.22e3}{K}}$.
\end{solution}
